\scanpage{30}
\uline{Cor} Supposons $u\in G_0$, \deleted{$\unsure{u\widehat c}$} et $u\widehat c$ de la forme $\widehat c'$
\uline{Alors} $u=\mathrm{id}$

En effet, $\deleted{\unsure{u\widehat c}}$ \uuline{$u.\widehat c=1.\widehat c'$} $\Rightarrow\ (u,c)=(1,c')$

Soit $G$ le groupe des autom \unsure{\deleted{qui}|qui} de $\widehat S$ qui sont l'identité mod $\mathcal J^3$ (donc $G/G_0\sim\Gamma(\Omega^{\otimes2})$) et considérons $(g,c)\longmapsto g.\widehat c$ sur $G\times\operatorname{Conn}(P)$. Quand a-t-on
\[
g\widehat c=g'\widehat c',\quad\text{i.e. }(g'^{-1}g)\widehat c=c'\quad\text{i.e. }h\widehat c=\widehat c'
\]
quand $h\widehat c$ est-il de la forme $\widehat c'$ ?

\inserted{Alors $c$ fixé, on trouve} \deleted{C'est} un sous-ensemble $G(c)$ de $G$
\[
h\in G(c)\quad\Longleftrightarrow\quad h.\widehat c\text{ de la forme }\widehat c'
\]
\deleted{$h\in G(c),\ g\in G_0,\ gh\in G(c)$}
On trouve que
\deleted{$\Rightarrow g=1$, donc pour $G(c)$}
\deleted{$G(c)\times$}
\diagram{\begin{tikzcd}
G_0\times G(c) \arrow[r] & G\\
(g,h) \arrow[r,maps to] & g.h
\end{tikzcd}}
est bijectif,

(si $gh=g'h'$ i.e. $\underbrace{(g'^{-1}g)}_{g''}h=h'$

d'où $\deleted{(g'^{-1}g)}\ (g''h).\widehat c$ de la forme $\widehat c''$
\[
\begin{array}{c}
\Vert\\
g''(\underbrace{h\widehat c}_{\widehat c'})
\end{array}
\]
D'où $g''=1$, d'où $\deleted{\unsure{g''}}\ g=g'$ et par suite $h=h'$.

\deleted{Inversement,} Surjectif, car soit $g\in G$
\diagram{\begin{tikzcd}[column sep=small]
\text{on a }g\widehat c= & k\widehat c' & \text{d'où } k^{-1}g\unsure{=|\equiv}h\in G(c)\\
& G_0 \arrow[u] &
\end{tikzcd}}
d'où $g=kh$
% archive foliation: 29
